\documentclass[10pt,a4paper]{amsart}
\usepackage[margin=19mm]{geometry}
\usepackage{amsmath,amssymb,mathtools}
\usepackage[scr=rsfs,cal=euler]{mathalfa}
\usepackage{xfrac}
\usepackage[unicode,hidelinks,bookmarks=false]{hyperref}
\newcommand{\ie}{i.e.\ }
\newcommand{\cf}{cf.\ }
\newcommand{\resp}{respectively\ }
\newcommand{\Subsection}{\subsection}
\providecommand{\nobreakdash}{\nobreak\hbox{-}}
\setlength{\emergencystretch}{2em}
\multlinegap0pt
\usepackage{amssymb}
\usepackage[shortlabels]{enumitem}
\usepackage{algorithm}\usepackage{algpseudocode}
\newtheorem{lemma}{Lemma}
\newcommand{\Hcal}{\mathcal{H}}
\newcommand{\R}{\mathbb{R}}
\newcommand{\grad}{\nabla}
\newcommand{\ip}[2]{\left\langle #1,#2 \right\rangle}
\newcommand{\norm}[1]{\left\lVert #1 \right\rVert}
\title{Curvature Recycling Douglas-Rachford Splitting: Transported Quasi-Newton Models for Expensive Smooth Proximal Subproblems}
\author{Musa Maharramov}
\date{Source: arXiv:2607.22895v1 (CC BY 4.0)}
\begin{document}
\maketitle

\noindent\emph{Source and licence.} Curvature Recycling Douglas-Rachford Splitting: Transported Quasi-Newton Models for Expensive Smooth Proximal Subproblems. Version arXiv:2607.22895v1 (CC BY 4.0), distributed by arXiv under the Creative Commons Attribution 4.0 International licence, http://creativecommons.org/licenses/by/4.0/, as recorded in the arXiv licence metadata for this exact version and shown on the abstract page https://arxiv.org/abs/2607.22895v1. Full licence notice: \texttt{licenses/arxiv-2607.22895-CC-BY-4.0.txt}.\par\smallskip

\noindent\textit{Editorial scope.} Counted unit: the article's lemma lem:certified-inner-contraction (source0.tex lines 609--618). The article's complete original proof of it is delivered with it (source0.tex lines 620--644, one \texttt{proof} environment of 25 lines). No mathematics is written, repaired or re-derived: the statement and the proof are the article's own lines, in the article's own notation, and no retained line is substituted or re-flowed.  Retained uncounted as the source-local context the proof needs: lines 563--572; lines 2198--2217.  Nothing was omitted from any retained span, so the package carries no omission record.  No retained line contains a citation, so the wrapper carries no bibliography and no bibliography entry.  No retained line contains a figure environment, an image-inclusion command or drawing code, so no image is delivered and the package declares no asset dependency.  Editorial formatting only, in addition: an AMS \texttt{amsart} preamble that restates the article's own declarations and macros used by the retained lines, namely the environments \texttt{lemma} on separate counters, together with the standard packages those lines need.  The retained environments print in this document as Lemma 1 in order of appearance, whereas the article prints the same items with the numbering of its own full document.  The complete native source of the article as distributed in the arXiv e-print for this exact version is bundled unchanged as \texttt{sources/205980/source0.tex}.
\begin{equation}
\label{eq:inner-theta}
    \theta_G:=\frac{L-m}{L+m},
    \qquad
    \theta:=\max\{\theta_Q,\theta_G^b\},
    \qquad
    L_p:=\frac1m,
    \qquad
    c:=2\alpha,
\end{equation}

\begin{lemma}[sector inequality for strongly convex smooth gradients]
\label{lem:sector}
Let $f:\Hcal\to\R$ be $a$-strongly convex and $b$-smooth, with
$0<a\le b$.  For $u,v\in\Hcal$, set
$s=u-v$ and $y=\grad f(u)-\grad f(v)$.  Then
\begin{equation}
\label{eq:sector-ineq}
  \ip{s}{y}
  \ge
  \frac{ab}{a+b}\norm{s}^2
  +\frac{1}{a+b}\norm{y}^2.
\end{equation}
Equivalently,
\begin{equation}
\label{eq:sector-ball-form}
  \norm{y-\tfrac{a+b}{2}s}
  \le\frac{b-a}{2}\norm{s}.
\end{equation}
The gradient difference therefore lies in the closed sector bounded by the slopes $a$ and $b$.
\end{lemma}

\begin{lemma}[Certified inner contraction]
\label{lem:certified-inner-contraction}
Let $p_k=P_{\gamma f}z_k$.  One call of the certified inner routine satisfies
\begin{equation}
\label{eq:inner-contraction-result}
    \norm{x_k-p_k}
    \le\theta\norm{a_{k-1}-p_k},
\end{equation}
with $\theta$ from \eqref{eq:inner-theta}.
\end{lemma}

\begin{proof}
If the quasi-Newton candidate is accepted, strong convexity and smoothness give
\[
    \norm{\widehat x_k-p_k}
    \le\frac1m\norm{r_k(\widehat x_k)}
    \le\frac{\theta_Q}{L}\norm{r_k(a_{k-1})}
    \le\theta_Q\norm{a_{k-1}-p_k}.
\]
If the candidate is rejected, apply the sector inequality from Lemma~\ref{lem:sector} to $r_k=\grad\phi_k$.  For one step $u^+=u-2r_k(u)/(m+L)$,
\[
\begin{aligned}
    \norm{u^+-p_k}^2
    &=\norm{u-p_k}^2
      -\frac{4}{m+L}\ip{u-p_k}{r_k(u)}
      +\frac{4}{(m+L)^2}\norm{r_k(u)}^2\\
    &\le\left(\frac{L-m}{L+m}\right)^2\norm{u-p_k}^2.
\end{aligned}
\]
After $b$ steps the factor is $\theta_G^b$.  Taking the larger of the two cases proves \eqref{eq:inner-contraction-result}.  The descent lemma also gives
\[
    \phi_k(u^+)\le\phi_k(u)
       -\frac{2m}{(m+L)^2}\norm{r_k(u)}^2,
\]
so every fallback step has a computable strict-descent certificate unless $u=p_k$.
\end{proof}

\end{document}
